\documentclass[aspectratio=169]{beamer}
\usetheme{Madrid}
\usecolortheme{seagull}
\usepackage[T1]{fontenc}
\usepackage{lmodern}
\usepackage{amsmath}
\usepackage{booktabs}

\title{FerrumMD: A Gentle Introduction to Molecular Dynamics in Rust}
\subtitle{What molecular dynamics is, what FerrumMD does, and where Rust helps}
\author{FerrumMD project overview}
\date{\today}

\begin{document}

\begin{frame}
  \titlepage
\end{frame}

\begin{frame}{Who this talk is for}
\begin{itemize}
  \item People who are curious about scientific computing, Rust, or molecular simulation.
  \item No molecular dynamics (MD) background is assumed.
  \item You only need three ideas:
  \begin{enumerate}
    \item atoms and molecules can be modeled as particles,
    \item particles move because forces act on them,
    \item a computer simulation repeats a small update many times.
  \end{enumerate}
\end{itemize}
\end{frame}

\begin{frame}{What problem does molecular dynamics solve?}
\begin{itemize}
  \item Experiments can show what molecules do, but not always every tiny motion.
  \item MD creates a computational ``movie'' of atoms and molecules over time.
  \item Typical questions:
  \begin{itemize}
    \item How does a liquid arrange itself?
    \item Does a molecule remain stable?
    \item How do temperature, pressure, or force-field choices affect behavior?
  \end{itemize}
  \item FerrumMD is a Rust codebase for building and studying these simulation workflows.
\end{itemize}
\end{frame}

\begin{frame}{MD in one minute}
\begin{block}{Basic idea}
Start with positions and velocities, compute forces, then move everything a tiny bit.
\end{block}
\begin{enumerate}
  \item Read or build a molecular system.
  \item Compute forces from the current geometry.
  \item Update positions and velocities using Newton's laws.
  \item Keep particles inside the simulation box and apply controls.
  \item Save energies, temperatures, and trajectory frames.
\end{enumerate}
\vspace{0.5em}
\centering
\texttt{state $\rightarrow$ forces $\rightarrow$ updated state $\rightarrow$ repeat}
\end{frame}

\begin{frame}{A few useful words}
\begin{description}
  \item[Particle] An atom or coarse-grained bead with a position, velocity, and mass.
  \item[Force field] Parameters and formulas that approximate molecular forces.
  \item[Integrator] The numerical method that advances the simulation clock.
  \item[Trajectory] The saved sequence of molecular snapshots.
  \item[Ensemble] The thermodynamic conditions being approximated, such as constant energy or temperature.
\end{description}
\end{frame}

\begin{frame}{Why this project matters}
\begin{itemize}
  \item \textbf{Learning value:} FerrumMD exposes the pieces of an MD engine in a readable systems language.
  \item \textbf{Engineering value:} Rust's ownership model helps organize mutable simulation state safely.
  \item \textbf{Scientific value:} explicit units, tests, and reproducible workflows are central to trustworthy simulations.
  \item \textbf{Scope today:} the main building blocks rather than production-scale chemistry.
\end{itemize}
\end{frame}

\begin{frame}{FerrumMD at a glance}
\begin{columns}[T]
\column{0.48\textwidth}
\textbf{Simulation core}
\begin{itemize}
  \item particle and molecule state
  \item simulation box and periodic boundaries
  \item Lennard--Jones and bonded forces
  \item velocity-Verlet time stepping
\end{itemize}

\column{0.48\textwidth}
\textbf{Workflow support}
\begin{itemize}
  \item temperature and pressure controls
  \item input/output for common file formats
  \item example water-like systems
  \item optional viewer, Python, and MPI paths
\end{itemize}
\end{columns}
\end{frame}

\begin{frame}{The simulation state: what the program stores}
\begin{itemize}
  \item \textbf{Coordinates:} where each particle is in the box.
  \item \textbf{Velocities:} how fast each particle is moving.
  \item \textbf{Forces:} the current push or pull on each particle.
  \item \textbf{Parameters:} masses, charges, bond constants, and interaction strengths.
  \item \textbf{Box:} the finite volume that stands in for a larger material.
\end{itemize}
\begin{block}{Unit convention}
FerrumMD keeps units explicit: nm, ps, amu, and kJ/mol.
\end{block}
\end{frame}

\begin{frame}{Forces: how particles influence each other}
\begin{columns}[T]
\column{0.52\textwidth}
\textbf{Nonbonded interactions}
\begin{itemize}
  \item Lennard--Jones models short-range attraction and repulsion.
  \item Periodic boundaries make the small box behave more like bulk matter.
  \item Cell subdivision reduces the number of pair checks.
\end{itemize}

\column{0.45\textwidth}
\textbf{Bonded interactions}
\begin{itemize}
  \item Harmonic bonds act like springs between connected atoms.
  \item Small molecular examples make the math easier to inspect.
  \item Topology readers help import force-field definitions.
\end{itemize}
\end{columns}
\end{frame}

\begin{frame}{Time stepping: turning forces into motion}
\begin{itemize}
  \item FerrumMD uses velocity-Verlet as a central integration method.
  \item Each step is tiny, so many steps are needed for a useful trajectory.
  \item The integrator must balance speed, stability, and energy behavior.
  \item Initialization can sample velocities from a Maxwell--Boltzmann distribution.
\end{itemize}
\begin{block}{Beginner mental model}
It is like a physics game loop, but numerical accuracy and units matter much more.
\end{block}
\end{frame}

\begin{frame}{Temperature, pressure, and ensembles}
\begin{itemize}
  \item Real experiments often control temperature or pressure.
  \item MD approximates these conditions with \emph{ensembles}:
  \begin{itemize}
    \item NVE: constant particles, volume, and energy,
    \item NVT: constant particles, volume, and temperature,
    \item NPT: constant particles, pressure, and temperature.
  \end{itemize}
  \item FerrumMD includes thermostat and barostat paths, including Berendsen and Nose--Hoover thermostats.
\end{itemize}
\end{frame}

\begin{frame}{Inputs, outputs, and visualization}
\begin{itemize}
  \item Read starting structures and parameters from familiar molecular formats.
  \item Supported inputs include PDB, GRO, LAMMPS data/dump, Martini, and CHARMM-style files.
  \item Write GRO/XTC outputs so trajectories can be viewed in external tools such as VMD.
  \item Optional Rust viewer and Python bindings provide extension points for demos and analysis.
\end{itemize}
\end{frame}

\begin{frame}{The honest Rust pitch}
\begin{block}{Not the claim}
Rust does \emph{not} make an MD engine automatically faster than C++.
Well-written C++ MD engines can be extremely fast.
\end{block}
\begin{block}{The stronger claim}
Rust offers C/C++-level native performance while making many dangerous scientific-code bugs harder to write.
\end{block}
\begin{itemize}
  \item The goal is not to beat mature engines just by choosing Rust.
  \item The goal is a safe, modular, maintainable path from research prototype to serious simulation software.
\end{itemize}
\end{frame}

\begin{frame}{Rust advantage 1: memory safety without GC}
\begin{itemize}
  \item MD engines manipulate large arrays of positions, velocities, forces, neighbor lists, cell lists, bonds, and topology.
  \item C/C++ bugs can include dangling pointers, use-after-free, double free, out-of-bounds access, and invalidated references.
  \item Rust's ownership and borrowing rules catch many of these hazards at compile time.
  \item Example MD hazard: mutating a particle array while another part of the code still holds references into it.
\end{itemize}
\begin{block}{Why this matters}
Force loops, neighbor-list builds, integration, and boundary-condition updates are all mutation-heavy and aliasing-sensitive.
\end{block}
\end{frame}

\begin{frame}{Rust advantage 2: safer parallelism}
\begin{itemize}
  \item MD wants parallelism in pair forces, domain decomposition, neighbor-list builds, reductions, and trajectory analysis.
  \item The danger is two threads accidentally writing the same force entry or sharing mutable state unsafely.
  \item Rust forces shared mutation to be explicit: locks, atomics, chunked mutable slices, or safe parallel iterators.
  \item Good MD designs usually avoid ``mutex everywhere'':
  \begin{itemize}
    \item each worker owns a particle chunk, or
    \item each worker writes to its own force buffer, then buffers are reduced safely.
  \end{itemize}
\end{itemize}
\end{frame}

\begin{frame}{Rust advantage 3: encode simulation meaning in types}
\begin{itemize}
  \item An MD engine is not just math; it is a large state machine.
  \item State includes box dimensions, cutoffs, thermostat/barostat state, masses, topology, random-number generators, and step counters.
  \item Rust types can make invalid states harder to represent:
  \begin{itemize}
    \item \texttt{Positions}, \texttt{Velocities}, \texttt{Forces}, and \texttt{Masses} instead of anonymous arrays,
    \item enums for algorithm choices such as no thermostat, Andersen, Berendsen, or Nose--Hoover,
    \item separate immutable parameters from mutable evolving state.
  \end{itemize}
  \item The compiler becomes a strict reviewer when ownership or invariants change.
\end{itemize}
\end{frame}

\begin{frame}{Rust advantage 4: tooling and modular architecture}
\begin{columns}[T]
\column{0.48\textwidth}
\textbf{Cargo workflow}
\begin{itemize}
  \item \texttt{cargo build}
  \item \texttt{cargo test}
  \item \texttt{cargo bench}
  \item \texttt{cargo doc}
  \item \texttt{cargo fmt} and \texttt{cargo clippy}
\end{itemize}

\column{0.48\textwidth}
\textbf{Engine design}
\begin{itemize}
  \item traits for force fields, integrators, and thermostats
  \item feature flags for MPI, viewer, Python, or future GPU paths
  \item unit tests for forces and regression tests for energy behavior
  \item clean APIs instead of one giant procedural script
\end{itemize}
\end{columns}
\end{frame}

\begin{frame}{Rust advantage 5: FFI and research workflows}
\begin{itemize}
  \item Rust can interoperate with C, C++, Python, MPI libraries, and existing visualization or analysis tools.
  \item A practical architecture is:
  \begin{itemize}
    \item Rust core engine for simulation kernels,
    \item Python interface for setup, analysis, and experimentation,
    \item file-format compatibility for tools such as VMD.
  \end{itemize}
  \item This supports a research workflow: prototype quickly, test carefully, and keep the performance-critical core compiled and typed.
\end{itemize}
\end{frame}

\begin{frame}{Where C++ still has advantages}
\begin{itemize}
  \item C++ has a larger and older HPC ecosystem.
  \item CUDA/HIP and GPU tooling are more mature in C++ workflows.
  \item Many scientific developers already know C++.
  \item Existing engines such as GROMACS, LAMMPS, NAMD, AMBER, and OpenMM have decades of optimization and validation.
\end{itemize}
\begin{block}{Positioning FerrumMD}
FerrumMD should be pitched as a modern, safe, modular research engine---not as faster than mature production engines simply because it uses Rust.
\end{block}
\end{frame}

\begin{frame}{How to read the codebase as a newcomer}
\begin{enumerate}
  \item Start with a runnable example to see the input, loop, and output flow.
  \item Inspect the state and molecule modules to understand stored data.
  \item Follow one force calculation from inputs to resulting forces and energy.
  \item Then look at integrators, ensembles, and I/O formats.
\end{enumerate}
\begin{block}{Tip}
You do not need to understand every force-field detail before contributing to tooling, tests, examples, or documentation.
\end{block}
\end{frame}

\begin{frame}{Roadmap and open opportunities}
\begin{itemize}
  \item Rigid-water constraints such as SETTLE or SHAKE/RATTLE.
  \item Topology-correct nonbonded exclusions and special pairs.
  \item Electrostatics and units audits for stronger TIP3P fidelity.
  \item Rust-focused opportunities:
  \begin{itemize}
    \item clearer examples and validation tests,
    \item SIMD or GPU experiments,
    \item stronger parallel decomposition beyond demo MPI,
    \item property-based tests for physical invariants.
  \end{itemize}
\end{itemize}
\end{frame}

\begin{frame}{Takeaways}
\begin{enumerate}
  \item MD is a repeated loop: compute forces, update motion, save observations.
  \item FerrumMD provides an approachable Rust implementation of that loop and its surrounding workflow.
  \item New contributors can start with examples, units, I/O, tests, and documentation before tackling advanced physics.
\end{enumerate}
\vspace{0.8em}
\centering
\large Questions \& discussion
\end{frame}

\end{document}
